Eranskin honetan nodoa martxan jartzeko beharrezko aginduak eta parametro-egiturak zehazten dira.

\needspace{6\baselineskip}
\section{Exekuzio-aginduak}\label{sec:exekuzioa}
Simulazioa eta paretari jarraitzeko nodoa abiarazteko urratsak:
\begin{enumerate}[leftmargin=2em, itemsep=0.3em]
  \item \textbf{Gazebo simulatzailea abiarazi}:
\begin{verbatim}
ros2 launch gazebo_worlds gz_robot.launch.py world:=yahboom_track
\end{verbatim}
  \item \textbf{P kontrolagailua exekutatu} ($K_p = 1{,}20$):
\begin{verbatim}
ros2 launch wall_following wf_ls.launch.py kp:=1.2 ki:=0.0 kd:=0.0
\end{verbatim}
  \item \textbf{PID kontrolagailu osoa exekutatu} ($K_p = 1{,}10,\ K_i = 0{,}04,\ K_d = 0{,}12$):
\begin{verbatim}
ros2 launch wall_following wf_ls.launch.py kp:=1.1 ki:=0.04 kd:=0.12
\end{verbatim}
\end{enumerate}

\needspace{6\baselineskip}
\section{Paketearen fitxategi-egitura}\label{sec:fitxategiak}
Garatutako soluzioaren fitxategi nagusiak:
\begin{itemize}[leftmargin=2em, itemsep=0.3em]
  \item \texttt{wall\_follow\_ls.py}: Kontrolagailu nagusia, proiekzio polar-kartesiarra, erregresio lineala eta PID algoritmo osoa jasotzen dituena.
  \item \texttt{wf\_ls.launch.py}: Parametro dinamikoak (\texttt{kp}, \texttt{ki}, \texttt{kd}) komando-lerrotik pasatzeko konfigurazioa.
  \item \texttt{results/}: Exekuzio bakoitzaren amaieran erroreen bilakaera eta distantziaren telemetria biltzen dituzten grafikoak gordetzen diren direktorioa.
\end{itemize}
